\item En 1993 el gobierno estadounidense instituyó el requisito de que todos los acondicionadores de aire que se vendan en este país deben tener un factor de eficiencia energética (EER, por sus siglas en inglés) de 10 o mayor. El EER se define como la razón de la capacidad de enfriamiento del acondicionador de aire, medida en unidades térmicas británicas por hora, o Btu/h, a su requerimiento de potencia eléctrica en watts. 

(a) Convierta el EER de 10.0 a forma adimensional, con el uso de la conversión $1\text{ Btu} = 1055\text{ J}$. 
(b) ¿Cuál es el nombre apropiado para esta cantidad adimensional? 
(c) Alrededor de 1970 era común encontrar acondicionadores de aire de habitación con EER de 5 o menos. Establezca la comparación de los costos de operación para acondicionadores de aire de $10\,000\text{ Btu/h}$ con EER de 5.00 y 10.0. Suponga que cada acondicionador de aire opera durante $1\,500\text{ h}$, en el verano, en una ciudad donde la electricidad cuesta 17.0 centavos por kWh.
